\section{Architecture du Système}

\subsection{Vue d'Ensemble}

Le système Moroccan Legal RAG suit une architecture en couches organisée
autour de deux flux principaux : le flux d'ingestion (PDF $\rightarrow$ base
de connaissances vectorielle) et le flux d'inférence (question $\rightarrow$
réponse générée). L'architecture est conçue pour être modulaire, chaque
couche étant interchangeable indépendamment.

\begin{figure}[H]
\centering
\fbox{
\begin{minipage}{0.9\columnwidth}
\footnotesize\ttfamily
\begin{tabular}{l}
\textbf{Flux d'Ingestion :} \\
PDF $\xrightarrow{\text{extraction}}$ Texte brut $\xrightarrow{\text{nettoyage}}$ \\
$\xrightarrow{\text{normalisation}}$ Texte normalis\'e $\xrightarrow{\text{parsing}}$ \\
Articles JSON $\xrightarrow{\text{embedding}}$ Vecteurs 1024d $\xrightarrow{\text{insertion}}$ pgvector \\
\\
\textbf{Flux d'Inf\'erence :} \\
Question $\xrightarrow{\text{embedding}}$ Vecteur $\xrightarrow{\text{recherche}}$ \\
Dense (pgvector cos) + BM25 (rank-bm25) \\
$\xrightarrow{\text{RRF}}$ Articles $\xrightarrow{\text{prompt}}$ LLM $\rightarrow$ R\'eponse
\end{tabular}
\end{minipage}
}
\caption{Architecture g\'en\'erale du syst\`eme Moroccan Legal RAG}
\end{figure}

\subsection{Composants du Flux d'Ingestion}

Le flux d'ingestion transforme les documents juridiques (PDF) en une base de
connaissances vectorielle interrogable. Il comprend cinq étapes principales :

\subsubsection{Classification et Extraction PDF avec Correction VLM}

Un classifieur de PDF \cite{src:pdfclassifier} détermine si un document est
numérique (texte intégral extractible) ou scanné (basé sur image), en
analysant le ratio pages textuelles/pages image avec des seuils configurables
dans \texttt{config/config.yaml}. Pour les PDF numériques, un extracteur
séquentiel essaye d'abord pdfplumber puis PyMuPDF (fallback layout-aware avec
reconstruction de colonnes).

Pour les PDF scannés, le pipeline complet comprend quatre étapes :
\begin{enumerate}
    \item \textbf{Rendu} : PyMuPDF convertit chaque page en image PNG
    \item \textbf{OCR} : EasyOCR effectue la reconnaissance du texte arabe
    sur GPU
    \item \textbf{Correction par VLM} : Gemma 4 en mode vision corrige les
    erreurs de reconnaissance avec un prompt spécialisé (\textit{«Corrige
    UNIQUEMENT les erreurs d'OCR. NE TRADUIS PAS. NE RESUMES PAS.»})
    \item \textbf{Post-traitement} : correction de l'ordre visuel arabe
    (inversion des segments de droite à gauche)
\end{enumerate}

Cette approche en cascade (OCR $\rightarrow$ VLM) permet de corriger les
erreurs fréquentes d'OCR pour l'arabe : confusion entre caractères
visuellement similaires (\textarabic{ت} / \textarabic{ث} / \textarabic{ب}),
segmentation incorrecte des mots liés, et omission de points diacritiques.

\subsubsection{Nettoyage du Texte Juridique}

Le module \texttt{LegalTextCleaner} \cite{src:cleaner} élimine le bruit
spécifique aux documents juridiques marocains :

\begin{itemize}
    \item \textbf{En-têtes et pieds de page} : motifs répétitifs détectés par
    fréquence de ligne sur l'ensemble du document
    \item \textbf{Numéros de page} : motifs \texttt{-\textbackslash d+-} et
    \texttt{\textbackslash d+/\textbackslash d+}
    \item \textbf{Watermarks et tampons} : textes fantômes et filigranes
    \item \textbf{Table des matières} : lignes avec motifs
    \texttt{......} et numéros de page
    \item \textbf{Duplication} : paragraphes identiques consécutifs
    \item \textbf{Sauts de ligne incorrects} : lignes coupées au milieu d'une
    phrase (détection par absence de ponctuation finale et longueur < 80 car.)
    \item \textbf{Artéfacts non textuels} : séquences de caractères spéciaux,
    espaces multiples, lignes vides excessives
\end{itemize}

\subsubsection{Normalisation Arabe}

Le module \texttt{ArabicNormalizer} \cite{src:normalizer} effectue une
normalisation standardisée du texte arabe :

\begin{itemize}
    \item \textbf{Normalisation Unicode NFKC} : décomposition des caractères
    composés
    \item \textbf{Unification des alefs} : \texttt{آ}, \texttt{أ}, \texttt{إ},
    \texttt{ٱ} $\rightarrow$ \texttt{ا}
    \item \textbf{Suppression des diacritiques} : voyelles courtes
    (fatha, damma, kasra), shadda, sukun
    \item \textbf{Suppression du tatweel} : caractère de prolongation
    (kashida) \texttt{ـ}
    \item \textbf{Normalisation du ya et du taa marbouta} : \texttt{ى}
    $\rightarrow$ \texttt{ي}, \texttt{ة} $\rightarrow$ \texttt{ه} en
    position finale
    \item \textbf{Normalisation des chiffres} : conversion des chiffres
    arabo-indiens orientaux (\texttt{٠١٢٣...}) en chiffres arabes standards
    (\texttt{0123...})
    \item \textbf{Espacement de la ponctuation} : séparation des signes de
    ponctuation des mots adjacents
\end{itemize}

\subsubsection{Analyse Structurelle et Parsing}

Le \texttt{LegalStructureParser} \cite{src:structure} extrait la structure
hiérarchique des codes de loi. Le texte normalisé est parsé pour identifier :

\begin{itemize}
    \item \textbf{Livres / Parties} : titres de haut niveau (\textarabic{القسم الأول})
    \item \textbf{Chapitres} : (\textarabic{الباب الأول}, \textarabic{الفصل الأول})
    \item \textbf{Sections} : (\textarabic{المبحث الأول}, \textarabic{الفرع الأول})
    \item \textbf{Articles} : motifs \texttt{المادة $X$} avec leur
    contenu textuel complet
\end{itemize}

Chaque article est ensuite enrichi avec son contexte hiérarchique complet
(nom de la loi, numéro de loi, livre, chapitre, section) et sérialisé au
format JSON pour l'ingestion.

\subsubsection{Ingestion Vectorielle}

Le pipeline d'ingestion \cite{src:ingestion} traite les articles JSON comme
suit :

\begin{enumerate}
    \item \textbf{Contextualisation} : chaque article est préfixé par son
    contexte : \texttt{[Code P\'enal] [Section 3] \textarabic{المادة 147:} ...}
    \item \textbf{Embedding} : le modèle \texttt{BAAI/bge-m3} produit un
    vecteur de 1024 dimensions. Les articles longs (>2000 caractères) sont
    tronqués pour éviter les OOM GPU.
    \item \textbf{Insertion} : les vecteurs sont insérés dans la table
    \texttt{legal\_chunks} de PostgreSQL, avec un index HNSW (m=16,
    ef\_construction=200) pour la recherche approximative du plus proche
    voisin (ANN).
\end{enumerate}

\subsection{Composants du Flux d'Inférence}

\subsubsection{API REST (FastAPI)}

L'API \cite{src:api} expose six routes sous \texttt{/api/v1} :

\begin{itemize}
    \item \texttt{GET /health} — État du système
    \item \texttt{POST /search} — Recherche pure (retrieval sans LLM)
    \item \texttt{POST /ask} — Q\&R complète (retrieval + génération LLM)
    \item \texttt{GET /laws} — Liste des lois et articles disponibles
    \item \texttt{POST /upload} — Upload et ingestion de nouveaux PDF
    \item \texttt{GET /conversations} — Historique des sessions Q\&R
\end{itemize}

\subsubsection{Moteur de Recherche Hybride}

Le système implémente trois stratégies de recherche interchangeables via un
pattern Strategy \cite{src:retrieval} :

\textbf{Recherche dense (vectorielle)} : la requête est encodée en un vecteur
1024d via le même modèle d'embedding. Une requête SQL avec l'opérateur de
distance cosinus \texttt{<=>} de pgvector trouve les $k$ plus proches voisins,
optionnellement filtrée par domaine juridique.

\textbf{Recherche BM25 (lexicale)} : la recherche plein-texte PostgreSQL
avec le dictionnaire arabe (\texttt{to\_tsvector(`arabic', ...)}) a été
abandonnée en raison de performances insuffisantes. Un index BM25Okapi a
été implémenté en mémoire avec un tokenizer arabe spécialisé
\cite{src:tokenizer}. Le tokenizer applique : (1) la normalisation via
l'\texttt{ArabicNormalizer} du projet, (2) l'extraction de mots par
expression régulière Unicode (plage arabe), (3) la suppression de 60+ mots
vides juridiques, et (4) le filtrage des tokens courts ($< 2$ caractères).
Les scores BM25 sont calculés avec les paramètres par défaut ($k_1 = 1,5$,
$b = 0,75$). Notre benchmark (section 5) confirme que cette approche atteint
un Recall@5 de 41,67\% et un Hit@5 de 46,88\% sur 32 questions juridiques.

\textbf{Reciprocal Rank Fusion (RRF)} : les $N$ meilleurs résultats de chaque
stratégie ($N = 20$) sont fusionnés selon la formule :

\[
\text{RRF}(d) = \sum_{s \in \{dense, bm25\}} \frac{1}{k + r_s(d)}
\]

où $k = 60$ est la constante de fusion et $r_s(d)$ est le rang du document
$d$ dans la liste de la stratégie $s$. Ce paramètre $k = 60$ a été choisi
empiriquement pour équilibrer l'influence des deux stratégies.

\subsubsection{Module de Génération (LLM)}

Le module de génération \cite{src:generation} construit un prompt structuré
à partir des articles récupérés et appelle un LLM distant. Le prompt inclut :

\begin{enumerate}
    \item Les \textbf{textes juridiques} récupérés, numérotés et sourcés
    \item Une \textbf{instruction système} en arabe spécifiant le rôle
    (assistant juridique marocain), la règle de grounding (répondre
    uniquement à partir des textes fournis), et le format de sortie
    \item La \textbf{question} de l'utilisateur
\end{enumerate}

Le LLM utilisé est Gemma 4 2B \cite{google2025gemma4}, quantifié en
Q4\_K\_XL et servi via llama.cpp \cite{gerganov2023llamacpp} sur un GPU
NVIDIA T4 (16 Go VRAM) hébergé sur Google Colab, accessible via un tunnel
Cloudflare. Le service LLM expose une API REST (\texttt{/generate},
\texttt{/health}) et gère le chargement du modèle, l'inférence, et le
streaming. La température est fixée à 0,3 pour favoriser la reproductibilité
des réponses, avec \texttt{max\_tokens}=4096 et \texttt{top\_p}=0,9.

La réponse générée est parsée en trois sections :
\begin{itemize}
    \item \textbf{Réponse courte} : résumé direct (1-2 phrases)
    \item \textbf{Analyse juridique} : explication détaillée avec citations
    \item \textbf{Exemple pratique} : application concrète (optionnel)
\end{itemize}

Un disclaimer automatique est ajouté à chaque réponse précisant la nature
automatique de la génération.

\subsection{Infrastructure de Données}

La base de données PostgreSQL 16 avec l'extension pgvector 0.7.4 stocke :

\begin{itemize}
    \item \texttt{legal\_chunks} : articles avec embedding vector(1024),
    métadonnées (domaine, loi, article, chapitre, source), et index HNSW
    \item \texttt{conversations} : historique Q\&R (session, rôle,
    contenu, sources JSONB)
    \item \texttt{app\_config} : configuration dynamique
\end{itemize}

\subsection{Interface Utilisateur}

L'interface web \cite{src:web} est développée avec SvelteKit 5 et Tailwind
CSS 4, avec support complet de l'arabe (RTL). Elle offre :

\begin{itemize}
    \item Page d'accueil avec champ de recherche juridique
    \item Page de résultats avec réponses structurées et sources citées
    \item Explorateur de lois avec arborescence hiérarchique
    \item Historique des sessions Q\&R
    \item Interface d'upload et de gestion des PDF
    \item Tableau de bord de santé du système
\end{itemize}

L'application supporte le multilingue (arabe, français) avec un système
d'internationalisation intégré.